\documentclass[12pt,a4paper]{article}

\usepackage[utf8]{inputenc}
\usepackage[T1]{fontenc}
\usepackage{lmodern}
\usepackage[margin=1in]{geometry}
\usepackage{setspace}
\usepackage{parskip}
\usepackage{graphicx}
\usepackage{enumitem}
\usepackage{cite}
\usepackage{xcolor}
\usepackage{amssymb}
\usepackage{booktabs}
\usepackage{tabularx}
\usepackage[font=small,labelfont=bf]{caption}
\usepackage[hyphens]{url}
\usepackage[colorlinks=true,linkcolor=black,citecolor=black,urlcolor=blue]{hyperref}

\onehalfspacing
\setlist{itemsep=0.6em}

% Sections are numbered 2.1, 2.2, ... beneath the unnumbered chapter headings.
\renewcommand{\thesection}{2.\arabic{section}}
\renewcommand{\thesubsection}{\thesection.\arabic{subsection}}
\renewcommand{\thetable}{2.\arabic{table}}
\renewcommand{\thefigure}{2.\arabic{figure}}

\newcommand{\cellyes}{$\checkmark$}
\newcommand{\cellpart}{$\circ$}
\newcommand{\cellno}{--}

\begin{document}

\begin{center}
  {\LARGE\bfseries Rule-Constrained Reinforcement Learning\\for Ambiguity-Aware IT Incident Triage\par}
  \vspace{1.5em}
  {\large OPOL CHARIS\\
  S24B23/003\\
  B29654\par}
\end{center}
\vspace{1.5em}

\section*{CHAPTER TWO}

\section*{2. LITERATURE REVIEW}

This chapter reviews the literature that bears on the five research questions of the study. It first describes how the literature was gathered and organised. It then synthesises the literature under five themes, one for each research question, and in each theme it separates what the reviewed studies establish from what they leave open. The chapter closes by comparing the studies closest to this one and by stating the research gap that the study addresses.

\section{Review Approach}

The review draws on three bodies of literature that Chapter One identified as relevant: AI-assisted IT incident management, safe and constrained reinforcement learning, and the deferral of decisions from automated systems to humans. A core set of 34 publications was collected across these areas. It comprises surveys of AIOps, empirical studies of incident management, foundational and recent work on shielding and constrained reinforcement learning, studies of selective prediction and learning to defer, and recent work on the runtime governance of language-model agents. Four further sources supply the standard formulations of reinforcement learning, partial observability, costly information acquisition and sequential deferral on which the discussion relies~\cite{sutton2018,kaelbling1998,janisch2019,joshi2023}.

A database search supplemented this core set. A search string combining three groups of terms, for the learning approach, for uncertainty and oversight, and for the incident-management setting, was run on Scopus, IEEE Xplore, the ACM Digital Library, SpringerLink and SIAM. The search returned 10,694 records with 10,541 distinct titles. These records were filtered by keyword for terms that denote IT operations and incident management. Seventy-seven records contained such terms in the title or abstract, and 19 of them referred specifically to incident management, incident response or AIOps. The abstracts of these records were read, and eight studies that addressed sequential decisions, information gathering, constraints or human involvement in incident handling were added to the review~\cite{alauthman2026,hammar2025,hammar2026llm,hammar2026belief,kim2025,dantonio2025,wang2023hrlhf,ding2023}. These eight studies are discussed on the basis of their abstracts, and the chapter attributes to them only what their abstracts state.

The synthesis is thematic. Each study was read for what it contributes to one or more of the research questions, and the findings were grouped into five themes that correspond to those questions: the representation of ambiguous incidents, sequential decisions about investigating, recommending and escalating, rule-based safety mechanisms, evaluation measures and baselines, and robustness under ambiguity. The review is therefore a structured thematic review of the collected literature. It does not claim the exhaustive coverage of a completed systematic review, and its conclusions about gaps are stated with that limit in mind.

\section{Representing Ambiguous Incident Triage in a Controlled Environment}

The first research question asks how ambiguous triage scenarios can be represented in a controlled environment. The literature answers part of this question by describing the evidence on which triage decisions rest. Remil et al.\ catalogue the data sources of incident management as incident reports, alerting signals, logs, metrics, traces and topology~\cite{remil2024}, and Zhang et al.\ add that language models have brought human-written material such as incident reports and documentation into use~\cite{zhang2025}. Incident reports are the primary source for the assignment of incidents~\cite{remil2024}. An environment for triage must therefore represent two kinds of evidence: the initial report, and the further observations that a technician could obtain from logs, metrics and other diagnostic sources.

The literature also explains why this evidence is ambiguous. Chen et al.\ found that distinct problems can produce similar symptoms, so that the symptoms alone do not identify the responsible team, and that incomplete knowledge of dependencies and imprecise assessment of resource health delay diagnosis~\cite{chen2020}. Remil et al.\ describe real IT data as commonly noisy and incomplete~\cite{remil2024}. Wu and Gong describe a related condition for agents that act in organisations: the facts needed to judge whether an action complies with policy may be absent from what the agent can see~\cite{wu2026}. Ambiguity in these accounts has three sources, which are missing observations, unreliable observations, and observations that are consistent with more than one cause. A controlled environment that claims to represent ambiguity should be able to produce each of them.

Realistic incident data is difficult to obtain, and this limits how the reviewed studies represent the task. Chen et al.\ report their results as relative values because of company policy~\cite{chen2020}. Remil et al.\ observe that publishing datasets raises confidentiality concerns, and that most methods for assignment and deduplication have been developed on software bug reports and not on operational incidents~\cite{remil2024}. Studies of sequential decision-making have responded by building or borrowing controlled environments. Dunsin et al.\ built a laboratory of infected virtual machines~\cite{dunsin2025}, Hammar and Li evaluated on the CAGE-2 benchmark~\cite{hammar2025}, Kim and D\'{a}n used simulations on synthetic and real dependency graphs~\cite{kim2025}, and Wang et al.\ combined production data with synthetic datasets containing injected anomalies~\cite{wang2023hrlhf}. Ren et al., in contrast, trained on a fixed dataset of recorded security incidents~\cite{ren2025}. Simulation is thus an accepted response to the scarcity of data, but the available environments model cyberattacks and microservice faults and not the triage of IT service incidents.

The formal tools for representing uncertainty about an incident are well established. A Markov decision process describes a decision problem in which the agent observes the state~\cite{sutton2018}, and a partially observable Markov decision process extends it to the case in which the agent receives only observations and must maintain a belief about the state~\cite{kaelbling1998}. Several incident-response studies adopt this view. Hammar et al.\ estimate a belief with particle filtering~\cite{hammar2026belief}, Kim and D\'{a}n maintain a belief in a hidden Markov model~\cite{kim2025}, and Alauthman et al.\ augment the state with recent telemetry and alert history to mitigate partial observability~\cite{alauthman2026}. These studies show that a hidden cause with noisy observations is a workable representation of an incident.

Taken together, the literature supplies the ingredients for representing ambiguous triage but not the representation itself. It identifies the evidence sources, explains the origins of ambiguity, and offers a formal model for hidden causes. What the reviewed studies do not provide is an environment for IT service incidents in which the level of ambiguity can be set deliberately and in which the available actions are those of triage. The first objective of this study responds to that absence.

\section{Reinforcement Learning for Investigating, Recommending and Escalating}

The second research question concerns how reinforcement learning can support the choice among gathering information, recommending an action and escalating. Research on incident triage has so far framed the task mainly as supervised prediction. The methods surveyed by Remil et al.\ range from Naive Bayes classifiers over report text to convolutional and recurrent networks that recommend a responsible team~\cite{remil2024}, and the triage component described by Chen et al.\ is a recurrent model over the incident text and discussion~\cite{chen2020}. One line of this work treats triage as a continuing process and updates its prediction as engineers add to the discussion~\cite{remil2024,chen2020}. The model in that case benefits from new evidence but does not decide to seek it. Zhang et al.\ add that automation in this field mainly stops at triage or the recommendation of solutions~\cite{zhang2025}. Where agents have been placed in the incident workflow, as in the runbook assistant that D'Antonio and Xie studied over 232 production incidents, the agent guides a human operator who carries out the mitigation~\cite{dantonio2025}. The triage literature therefore predicts well-defined labels but does not model the decision about what to do when the label is uncertain.

Reinforcement learning has entered IT operations and incident handling, mostly for diagnosis and response. Ding et al.\ learn a policy that prunes a service dependency graph before root cause analysis~\cite{ding2023}, and Wang et al.\ use hierarchical reinforcement learning to reduce the number of questions put to engineers when that graph is constructed~\cite{wang2023hrlhf}. In cybersecurity, Ren et al.\ train an agent to select response strategies~\cite{ren2025}, Dunsin et al.\ train an agent to conduct a malware investigation~\cite{dunsin2025}, and Alauthman et al.\ propose a deep Q-network agent for incident response in industrial control systems~\cite{alauthman2026}. Hammar and colleagues treat incident response as planning under partial observability~\cite{hammar2025,hammar2026belief}. These studies establish that sequential formulations of incident handling are feasible and can outperform fixed procedures in their own settings. Their action sets, however, consist of responses or investigative steps, and the studies do not address triage as defined in this study.

The decision to gather information has been studied as a sequential problem in its own right. Under partial observability, an action may be chosen for the information it yields, and its value depends on what that information costs~\cite{kaelbling1998}. Janisch et al.\ trained an agent to acquire costly features one at a time before classifying~\cite{janisch2019}. In incident handling, Kim and D\'{a}n choose which alerts to investigate so as to reduce the uncertainty of a belief~\cite{kim2025}, the agent of Dunsin et al.\ collects forensic evidence step by step~\cite{dunsin2025}, and Wang et al.\ learn which questions to ask engineers~\cite{wang2023hrlhf}. Information gathering is consequently a tractable learning problem, and one that has already been applied to parts of incident handling. In these studies the alternative to gathering information is to conclude, and none of them offers the agent the further option of passing the case to a person.

The decision to hand a case to a human has been studied separately, and mostly as a single-step choice. Selective prediction allows a model to abstain when its confidence is low~\cite{geifman2019}, and learning to defer trains a model to pass a case to a human decision-maker by taking that person's ability into account~\cite{madras2018,hemmer2023}. Strong et al.\ note that the standard formulation treats cases as independent and optimises a one-step outcome~\cite{strong2025}. Sequential extensions exist. Joshi et al.\ place deferral in a Markov decision process, where deferring affects later states~\cite{joshi2023}, and Zhang et al.\ formulate deferral as a constrained Markov decision process that accounts for the fatigue of the human expert~\cite{zhang2026falcon}. Guidance on AI risk management likewise lists deferral to a human expert among the configurations in which an AI system may operate~\cite{nist2023}. Applied work treats escalation as an estimate of when a model is likely to be wrong~\cite{bachar2026} or as a decision that policy can make mandatory~\cite{pacella2026}. In each case the choice is between acting and deferring, and the option of first gathering more evidence is absent.

The two lines of work thus stop on either side of the choice this study examines. Information-gathering studies offer the agent two options, to look further or to decide, and deferral studies offer two different options, to decide or to hand over. Hammar et al.\ come closest to joining them: their language-model planner abstains when its candidate actions are inconsistent and then collects external feedback, for example from a digital twin~\cite{hammar2026llm}. That system gathers feedback from a model of the system and not from a human, and its decision rule is a consistency threshold and not a learned policy. Among the reviewed studies, none trains a single policy to choose among investigating, recommending and escalating in IT incident triage.

\section{Rule-Based Safety Mechanisms and Operational Constraints}

The third research question asks how effectively a rule-based mechanism can prevent actions that violate operational constraints. Safe reinforcement learning offers two broad ways to impose constraints on a learning agent. Garc\'{\i}a and Fern\'{a}ndez divide the field into methods that modify the optimality criterion and methods that modify the exploration process with external knowledge~\cite{garcia2015}. Constrained policy optimisation belongs to the first group and bounds the expected cumulative cost of a policy~\cite{achiam2017,wachi2024}. Zhao et al.\ point out that many practical constraints are instead instantaneous and deterministic and must hold in every state~\cite{zhao2023}. An operational rule that forbids a particular action under stated conditions is a constraint of this second kind, which an expected-cost bound does not enforce.

Shielding enforces constraints of that kind by separating the rule from the learner. Wen et al.\ first computed a permissive strategy that satisfies a temporal logic specification and then let the agent learn within it~\cite{wen2015}. Alshiekh et al.\ generalised the idea into a shield that is synthesised from a specification and an abstraction of the environment, and that either restricts the actions offered to the agent or replaces an unsafe action after the agent selects it~\cite{alshiekh2018}. Later work relaxed the requirement of absolute safety. Jansen et al.\ block an action only when it raises the probability of a violation beyond a threshold~\cite{jansen2020}, and Yang et al.\ use probabilistic logic to handle noisy observations~\cite{yang2023}. Hamel-De le Court and colleagues have addressed scalability and uncertainty in the environment model~\cite{hamel2025,hamel2026}. K\"{o}nighofer et al.\ summarise the benefits of shielding as assured compliance with the specification, permissiveness towards safe actions, and preserved convergence of the learner~\cite{konighofer2025}.

Shielding has costs that matter for this study. A shield requires a formal model of the safety-relevant part of the environment, and its guarantee extends only as far as that model and the specification are correct~\cite{konighofer2025}. Pranger and K\"{o}nighofer observe that shielding has seen limited use in practice because it demands expertise in formal methods and considerable engineering effort~\cite{pranger2026}. Yang et al.\ note that a shield which simply rejects actions can hinder learning, because the agent is unaware that its choices were overridden~\cite{yang2023}. The demonstrations in this literature are grid worlds, arcade games, simulated robots and control tasks~\cite{alshiekh2018,jansen2020,achiam2017}. The method is therefore well understood in principle, while its use on the rule sets of IT operations remains largely unexamined.

A parallel body of work enforces explicit rules on language-model agents at run time. Wang et al.\ define a rule language with triggers, predicates and enforcement actions~\cite{wang2025agentspec}, Xiang et al.\ and Chen et al.\ use a second agent to check the actions of the first against stated policies~\cite{xiang2025,chen2025shieldagent}, and Pacella et al.\ place a policy-derived guardrail beneath the agent so that it cannot fall below the escalation that policy requires~\cite{pacella2026}. Wu and Gong check each proposed action against organisational state and return one of three verdicts, to allow, to block or to ask for clarification~\cite{wu2026}. Hu et al.\ show the scale that rule sets can reach, with a dataset of 12,985 rules drawn from 74 regulations and policies~\cite{hu2026}. This work confirms that rules kept outside the decision-maker can be inspected and enforced independently. It concerns language-model agents and not reinforcement learning, and it does not study how enforcement interacts with a policy that is still learning.

Two studies bring rule enforcement into incident response itself. Alauthman et al.\ couple their deep Q-network agent with an explicit safety shield that enforces operational constraints, in a design where the agent remains advisory by default and the operator can override it~\cite{alauthman2026}. Hammar et al.\ check the candidate actions of a language model for consistency with system constraints~\cite{hammar2026llm}. The first of these is the closest precedent for the mechanism proposed in this study, and it shows that the combination of a learned policy and a rule-based shield has been proposed for incident response. Its setting is cyber-physical control systems, its agent selects responses and does not choose whether to investigate or escalate, and its abstract reports an evaluation plan with indicative findings.

The literature therefore supports the feasibility of an independent rule-based mechanism and identifies the questions that remain. It is established that explicit rules can be enforced on a learning agent without preventing it from learning. It is less clear what should replace a blocked proposal when the agent's alternatives include requesting information and escalating, and how often a learned triage policy would propose prohibited actions in the first place. The third objective of this study addresses these questions for IT incident triage.

\section{Evaluation Measures and Baseline Policies}

The fourth research question concerns how the proposed approach should be compared with baseline policies. Each strand of the literature has its own measures. AIOps research evaluates triage with classification and ranking metrics and evaluates the process with time-based measures such as mean time to repair and mean time to engage~\cite{remil2024}. Remil et al.\ criticise the reliance on contingency-table metrics that ignore context and time~\cite{remil2024}, and Zhang et al.\ report that generated diagnoses and mitigation suggestions often require human assessment because they lack a fixed ground truth~\cite{zhang2025}. Reinforcement learning studies of incident response report the effectiveness of the response, the time taken, and false positive rates~\cite{ren2025}, or the time and accuracy of an investigation~\cite{dunsin2025}. These measures describe whether the final decision was correct and how quickly it was reached.

Studies of constraints and of deferral measure different things. Shielding studies compare a shielded agent with the same agent unshielded and report safety violations alongside reward~\cite{alshiekh2018,jansen2020}. Research on language-model agents has introduced measures that combine the two: Levy et al.\ credit a task only when it is completed without violating any applicable policy, and they find that this measure falls well below the raw completion rate for the agents they test~\cite{levy2026}. Mohammadi et al.\ list compliance with policy among the objectives that agent evaluation in enterprises must cover~\cite{mohammadi2025}. Deferral studies report accuracy as a function of coverage~\cite{geifman2019}, the accuracy of the combined human and model~\cite{madras2018}, or the trade-off between accuracy and the cost of human review~\cite{bachar2026}. Pacella et al.\ measure the recall of mandatory escalations and the stability of the decision across runs~\cite{pacella2026}. Turan relates the escalation rate to realised safety and argues that the relationship need not be monotonic~\cite{turan2026}.

The baselines used in these studies follow a common pattern. Learned policies are compared with fixed rule-based procedures~\cite{ren2025,pacella2026}, constrained agents with their unconstrained counterparts~\cite{alshiekh2018,jansen2020,pacella2026}, and deferring systems with the model alone and the human alone~\cite{madras2018,zhang2026falcon}. Dunsin et al.\ compare their agent with human experts~\cite{dunsin2025}. These comparisons isolate one property at a time, which is the contribution of learning, of the constraint, or of deferral.

No reviewed study reports all four kinds of measure for one agent. Decision quality, the cost of investigation, escalation behaviour and constraint compliance are each measured somewhere in the literature, but in different studies and for different systems. An agent that can investigate, recommend and escalate could score well on any one of them by neglecting the others, for example by escalating every case or by investigating without limit. The literature thus provides the individual measures and the baseline patterns that this study needs, and it leaves open how they behave when reported together. The fourth objective takes up that question.

\section{Robustness under Ambiguity and Missing Evidence}

The fifth research question asks how robust the approach is as evidence becomes incomplete, conflicting or more ambiguous. The literature treats robustness to imperfect data as a requirement. Remil et al.\ list it among the properties that AIOps solutions must have~\cite{remil2024}, and Mohammadi et al.\ treat robustness to variation as part of the reliability that agent evaluation should probe~\cite{mohammadi2025}. Some studies test it directly. Kim and D\'{a}n report that their investigation policies remain effective under high false-alert rates and investigation errors~\cite{kim2025}, and Hammar and Li evaluate response planning when the assumed model of the incident is wrong~\cite{hammar2025}. Ren et al.\ acknowledge that their agent generalises less well to attack patterns unlike those in its training data~\cite{ren2025}. These results concern the quality of the final decision as conditions degrade.

Uncertainty also affects the mechanisms that are meant to keep an agent within bounds. A deterministic shield assumes that the agent's observations are accurate, an assumption Yang et al.\ describe as unrealistic~\cite{yang2023}, and a shield computed from an estimated model can lose its guarantee when the estimate is wrong~\cite{hamel2026}. Wu and Gong show that a rule cannot be applied correctly when the facts it depends on are hidden from the checker~\cite{wu2026}. A rule-based mechanism in triage faces the same difficulty, because a rule whose condition refers to the state of a system can only be evaluated on the evidence available.

The human side of escalation is not constant either. Learning when to defer depends on data about the expert that is costly to obtain~\cite{hemmer2023}, people decide differently when they know that a model has deferred to them~\cite{bondi2022}, and expert accuracy changes with workload~\cite{zhang2026falcon}. Turan argues from a modelling study that increasing the escalation rate can reduce realised safety once the reviewer is overloaded~\cite{turan2026}. An agent that responds to ambiguity only by escalating more would therefore transfer the problem and not resolve it.

The reviewed studies have not examined how the balance among investigating, recommending and escalating shifts as ambiguity rises. Robustness has been measured for the accuracy of a decision, for the validity of a shield, and for the reliability of a human reviewer, each in isolation. For a triage agent, the behaviour of interest is the change in its mix of decisions: whether it investigates more, escalates more, or continues to recommend with less support. The fifth objective of this study measures that change across controlled levels of ambiguity and missing evidence.

\section{Synthesis of the Reviewed Literature}

Across the five themes the literature shows a consistent pattern: the components that this study needs exist, but they have been developed and evaluated apart. AIOps research describes the triage task and its evidence but models it as prediction. Reinforcement learning research on incident handling models sequential decisions but concentrates on response and investigation in cybersecurity. Shielding research provides a way to enforce rules on a learning agent but has been demonstrated mainly on control tasks. Deferral research explains when a model should hand a decision to a person but treats the hand-over mostly as a single-step alternative to prediction.

Table~\ref{tab:comparison} compares the studies closest to this one against the capabilities that the research questions concern. The comparison shows that each capability appears in at least one study and that the studies cluster in pairs of capabilities. Sequential policies are combined with information gathering~\cite{dunsin2025,kim2025,wang2023hrlhf}, or with a rule check~\cite{alauthman2026}, or with deferral~\cite{zhang2026falcon}. Rule checks are combined with escalation where the decision-maker is a language model~\cite{pacella2026}. No row other than the proposed study combines a sequential policy, information gathering, a rule check and human escalation, and none of the studies of IT service incidents includes a rule check or an escalation decision.

\begin{table}[htbp]
\centering
\caption{Comparison of the studies closest to this research. \cellyes~=~addressed; \cellpart~=~partly addressed; \cellno~=~not addressed or not reported. Studies marked $\dagger$ were assessed from their abstracts.}
\label{tab:comparison}
\small
\setlength{\tabcolsep}{4pt}
\begin{tabularx}{\textwidth}{@{}l>{\raggedright\arraybackslash}Xccccc@{}}
\toprule
Study & Setting & \rotatebox{70}{Sequential policy} & \rotatebox{70}{Information gathering} & \rotatebox{70}{Rule check} & \rotatebox{70}{Human escalation} & \rotatebox{70}{Varied ambiguity} \\
\midrule
Chen et al.~\cite{chen2020} & IT incident triage & \cellno & \cellno & \cellno & \cellno & \cellno \\
Wang et al.~\cite{wang2023hrlhf}$^\dagger$ & Root cause analysis & \cellyes & \cellyes & \cellno & \cellpart & \cellno \\
Ren et al.~\cite{ren2025} & Security incident response & \cellyes & \cellno & \cellno & \cellno & \cellno \\
Dunsin et al.~\cite{dunsin2025} & Malware investigation & \cellyes & \cellyes & \cellno & \cellno & \cellno \\
Kim and D\'{a}n~\cite{kim2025}$^\dagger$ & Security alert prioritisation & \cellyes & \cellyes & \cellno & \cellno & \cellyes \\
Hammar and Li~\cite{hammar2025}$^\dagger$ & Security incident response & \cellyes & \cellno & \cellno & \cellno & \cellyes \\
Hammar et al.~\cite{hammar2026llm}$^\dagger$ & Security response planning & \cellpart & \cellpart & \cellyes & \cellno & \cellno \\
Alauthman et al.~\cite{alauthman2026}$^\dagger$ & Industrial control incident response & \cellyes & \cellno & \cellyes & \cellpart & \cellno \\
Zhang et al.~\cite{zhang2026falcon} & Classification with deferral & \cellyes & \cellno & \cellno & \cellyes & \cellno \\
Pacella et al.~\cite{pacella2026} & Employee onboarding workflow & \cellno & \cellno & \cellyes & \cellyes & \cellno \\
Wu and Gong~\cite{wu2026} & Workplace assistant actions & \cellno & \cellpart & \cellyes & \cellno & \cellno \\
\midrule
This study (proposed) & IT incident triage & \cellyes & \cellyes & \cellyes & \cellyes & \cellyes \\
\bottomrule
\end{tabularx}
\end{table}

The comparison also qualifies any claim to novelty. A reinforcement learning agent with a safety shield has been proposed for incident response in industrial control systems~\cite{alauthman2026}, and a planner that checks its actions against constraints and seeks further feedback when uncertain has been evaluated for security incidents~\cite{hammar2026llm}. This study does not introduce the idea of constraining a learned incident-handling agent. Its contribution lies in the particular decision it studies, the setting in which it studies it, and the way it evaluates the result.

\section{Research Gap}

The review identifies five related gaps, one under each research question. The first is a gap in representation. The reviewed studies explain why triage evidence is ambiguous and provide formal models for hidden causes, but they offer no controlled environment for IT service incidents in which ambiguity can be varied and in which the actions are those of triage. The available environments model cyberattacks and microservice faults~\cite{hammar2025,dunsin2025,wang2023hrlhf}.

The second gap concerns the decision itself. Information-gathering studies let an agent choose between looking further and deciding~\cite{janisch2019,kim2025}, and deferral studies let it choose between deciding and handing over~\cite{madras2018,joshi2023,zhang2026falcon}. No reviewed study trains one policy to choose among all three options for IT incident triage. It is consequently unknown whether an agent can learn to treat investigation and escalation as distinct responses to uncertainty, or whether it will collapse into one of them.

The third gap concerns the interaction between rules and that decision. Rule-based shields have been shown to restrict learned policies in control tasks~\cite{alshiekh2018,konighofer2025}, and one study proposes a shield for incident response~\cite{alauthman2026}. The reviewed literature does not report how an independent rule-based mechanism behaves when the constrained agent can also request information or escalate, nor how often a learned triage policy proposes actions that the rules prohibit.

The fourth and fifth gaps concern evidence. Decision quality, investigation cost, escalation behaviour and constraint compliance have each been measured, but not together for one agent, so the trade-offs among them have not been observed. Robustness has been tested for decision accuracy and for the validity of shields~\cite{kim2025,hammar2025,hamel2026}, but not for the mix of investigating, recommending and escalating as ambiguity increases.

This study addresses these gaps within stated limits. It will construct a simulated environment for ambiguity-aware IT incident triage, train a rule-constrained reinforcement learning agent to choose among gathering diagnostic information, recommending a permitted action and escalating to a human technician, and evaluate the agent against baseline decision policies on all four measures and across levels of ambiguity. The gaps are stated for the literature reviewed in this chapter. Related work may exist beyond it, and the claim made here is that the reviewed literature leaves these questions open, not that they have never been examined.

\section{Chapter Summary}

This chapter reviewed the literature under five themes that correspond to the research questions of the study. It showed that the literature explains the sources of ambiguity in triage and supplies formal models for it, that sequential decision methods have been applied to incident response and investigation but not to the three-way choice of triage, that rule-based mechanisms can constrain learning agents but have rarely been examined in IT operations, and that evaluation measures and robustness tests exist separately for each strand. A comparison of the closest studies showed that none combines a sequential policy, information gathering, a rule check and human escalation, and the chapter stated five gaps on that basis. Chapter Three presents the methodology by which the study addresses them.

\begin{thebibliography}{99}

\bibitem{sutton2018} R. S. Sutton and A. G. Barto, \emph{Reinforcement Learning: An Introduction}, 2nd ed. Cambridge, MA, USA: MIT Press, 2018.

\bibitem{kaelbling1998} L. P. Kaelbling, M. L. Littman, and A. R. Cassandra, ``Planning and acting in partially observable stochastic domains,'' \emph{Artificial Intelligence}, vol. 101, no. 1--2, pp. 99--134, 1998. [Online]. Available: \url{https://doi.org/10.1016/S0004-3702(98)00023-X}

\bibitem{janisch2019} J. Janisch, T. Pevn\'{y}, and V. Lis\'{y}, ``Classification with costly features using deep reinforcement learning,'' in \emph{Proc. AAAI Conf. Artif. Intell.}, vol. 33, no. 1, 2019, pp. 3959--3966. [Online]. Available: \url{https://doi.org/10.1609/aaai.v33i01.33013959}

\bibitem{joshi2023} S. Joshi, S. Parbhoo, and F. Doshi-Velez, ``Learning-to-defer for sequential medical decision-making under uncertainty,'' \emph{Transactions on Machine Learning Research}, 2023. [Online]. Available: \url{https://openreview.net/forum?id=0pn3KnbH5F}

\bibitem{alauthman2026} M. Alauthman, A. Almomani, A. Al-Qerem, G. Albesani, and M. Alkasassbeh, ``Explainable reinforcement learning for automated incident response in critical infrastructure,'' in \emph{Proc. 2nd Int. Conf. Comput. Intell. Approaches Appl. (ICCIAA)}, 2026, pp. 1--7. [Online]. Available: \url{https://doi.org/10.1109/ICCIAA68481.2026.11543556}

\bibitem{hammar2025} K. Hammar and T. Li, ``Online incident response planning under model misspecification through Bayesian learning and belief quantization,'' in \emph{Proc. 18th ACM Workshop Artif. Intell. Secur. (AISec)}, 2025, pp. 40--51. [Online]. Available: \url{https://doi.org/10.1145/3733799.3762965}

\bibitem{hammar2026llm} K. Hammar, T. Alpcan, and E. C. Lupu, ``Hallucination-resistant security planning with a large language model,'' in \emph{Proc. IEEE Netw. Oper. Manage. Symp. (NOMS)}, 2026, pp. 1--12. [Online]. Available: \url{https://doi.org/10.1109/NOMS69089.2026.11668410}

\bibitem{hammar2026belief} K. Hammar, Y. Li, T. Alpcan, E. C. Lupu, and D. P. Bertsekas, ``Adaptive network security policies via belief aggregation and rollout,'' \emph{IEEE Transactions on Information Forensics and Security}, vol. 21, pp. 6544--6559, 2026. [Online]. Available: \url{https://doi.org/10.1109/TIFS.2026.3709131}

\bibitem{kim2025} Y. Kim and G. D\'{a}n, ``Dynamic alert prioritization for real-time situational awareness: A hidden Markov model framework with active learning,'' \emph{IEEE Transactions on Dependable and Secure Computing}, vol. 22, no. 6, pp. 6913--6928, 2025. [Online]. Available: \url{https://doi.org/10.1109/TDSC.2025.3592127}

\bibitem{dantonio2025} R. D'Antonio and H. Xie, ``Human-in-the-loop runbook improvement with agentic support automation,'' in \emph{Proc. IEEE 7th Int. Conf. Cogn. Mach. Intell. (CogMI)}, 2025, pp. 341--351. [Online]. Available: \url{https://doi.org/10.1109/CogMI67134.2025.00046}

\bibitem{wang2023hrlhf} L. Wang \emph{et al.}, ``Root cause analysis for microservice systems via hierarchical reinforcement learning from human feedback,'' in \emph{Proc. 29th ACM SIGKDD Conf. Knowl. Discov. Data Min. (KDD)}, 2023, pp. 5116--5125. [Online]. Available: \url{https://doi.org/10.1145/3580305.3599934}

\bibitem{ding2023} R. Ding \emph{et al.}, ``TraceDiag: Adaptive, interpretable, and efficient root cause analysis on large-scale microservice systems,'' in \emph{Proc. 31st ACM Joint Eur. Softw. Eng. Conf. Symp. Found. Softw. Eng. (ESEC/FSE)}, 2023, pp. 1762--1773. [Online]. Available: \url{https://doi.org/10.1145/3611643.3613864}

\bibitem{remil2024} Y. Remil, A. Bendimerad, R. Mathonat, and M. Kaytoue, ``AIOps solutions for incident management: Technical guidelines and a comprehensive literature review,'' arXiv preprint arXiv:2404.01363, 2024. [Online]. Available: \url{https://arxiv.org/abs/2404.01363}

\bibitem{zhang2025} L. Zhang \emph{et al.}, ``A survey of AIOps in the era of large language models,'' \emph{ACM Computing Surveys}, vol. 58, no. 2, Art. no. 44, 2025. [Online]. Available: \url{https://doi.org/10.1145/3746635}

\bibitem{chen2020} Z. Chen \emph{et al.}, ``Towards intelligent incident management: Why we need it and how we make it,'' in \emph{Proc. 28th ACM Joint Meeting Eur. Softw. Eng. Conf. Symp. Found. Softw. Eng. (ESEC/FSE)}, 2020, pp. 1487--1497. [Online]. Available: \url{https://doi.org/10.1145/3368089.3417055}

\bibitem{wu2026} J. Wu and M. Gong, ``Policy-invisible violations in LLM-based agents,'' arXiv preprint arXiv:2604.12177, 2026. [Online]. Available: \url{https://arxiv.org/abs/2604.12177}

\bibitem{dunsin2025} D. Dunsin, M. C. Ghanem, K. Ouazzane, and V. Vassilev, ``Reinforcement learning for an efficient and effective malware investigation during cyber incident response,'' \emph{High-Confidence Computing}, vol. 5, no. 3, Art. no. 100299, 2025. [Online]. Available: \url{https://doi.org/10.1016/j.hcc.2025.100299}

\bibitem{ren2025} S. Ren, J. Jin, G. Niu, and Y. Liu, ``ARCS: Adaptive reinforcement learning framework for automated cybersecurity incident response strategy optimization,'' \emph{Applied Sciences}, vol. 15, no. 2, p. 951, 2025. [Online]. Available: \url{https://doi.org/10.3390/app15020951}

\bibitem{geifman2019} Y. Geifman and R. El-Yaniv, ``SelectiveNet: A deep neural network with an integrated reject option,'' in \emph{Proc. 36th Int. Conf. Mach. Learn. (ICML)}, PMLR vol. 97, 2019, pp. 2151--2159.

\bibitem{madras2018} D. Madras, T. Pitassi, and R. Zemel, ``Predict responsibly: Improving fairness and accuracy by learning to defer,'' in \emph{Advances in Neural Information Processing Systems}, vol. 31, 2018.

\bibitem{hemmer2023} P. Hemmer, L. Thede, M. V\"{o}ssing, J. Jakubik, and N. K\"{u}hl, ``Learning to defer with limited expert predictions,'' in \emph{Proc. AAAI Conf. Artif. Intell.}, vol. 37, no. 5, 2023, pp. 6002--6011. [Online]. Available: \url{https://doi.org/10.1609/aaai.v37i5.25742}

\bibitem{strong2025} J. Strong, E. Sun, H. Rogers, H. Higham, and J. A. Noble, ``Learning to defer: A survey,'' preprint, Dec. 2025.

\bibitem{zhang2026falcon} Z. Zhang, C. C. Nguyen, D. Rosewarne, K. Wells, and G. Carneiro, ``Fatigue-aware learning to defer via constrained optimisation,'' arXiv preprint arXiv:2604.00904, 2026. [Online]. Available: \url{https://arxiv.org/abs/2604.00904}

\bibitem{nist2023} National Institute of Standards and Technology, ``Artificial Intelligence Risk Management Framework (AI RMF 1.0),'' NIST AI 100-1, Gaithersburg, MD, USA, Jan. 2023. [Online]. Available: \url{https://doi.org/10.6028/NIST.AI.100-1}

\bibitem{bachar2026} O. Bachar \emph{et al.}, ``LLM performance predictors: Learning when to escalate in hybrid human-AI moderation systems,'' arXiv preprint arXiv:2601.07006, 2026. [Online]. Available: \url{https://arxiv.org/abs/2601.07006}

\bibitem{pacella2026} M. Pacella, G. Papadia, and V. Giliberti, ``Governed agentic process automation: A floor-safety guarantee for compliance-critical LLM routing,'' \emph{Algorithms}, vol. 19, no. 8, p. 627, 2026. [Online]. Available: \url{https://doi.org/10.3390/a19080627}

\bibitem{garcia2015} J. Garc\'{\i}a and F. Fern\'{a}ndez, ``A comprehensive survey on safe reinforcement learning,'' \emph{Journal of Machine Learning Research}, vol. 16, no. 42, pp. 1437--1480, 2015.

\bibitem{achiam2017} J. Achiam, D. Held, A. Tamar, and P. Abbeel, ``Constrained policy optimization,'' in \emph{Proc. 34th Int. Conf. Mach. Learn. (ICML)}, PMLR vol. 70, 2017, pp. 22--31.

\bibitem{wachi2024} A. Wachi, X. Shen, and Y. Sui, ``A survey of constraint formulations in safe reinforcement learning,'' in \emph{Proc. 33rd Int. Joint Conf. Artif. Intell. (IJCAI)}, 2024, pp. 8262--8271. [Online]. Available: \url{https://doi.org/10.24963/ijcai.2024/913}

\bibitem{zhao2023} W. Zhao, T. He, R. Chen, T. Wei, and C. Liu, ``State-wise safe reinforcement learning: A survey,'' in \emph{Proc. 32nd Int. Joint Conf. Artif. Intell. (IJCAI)}, 2023, pp. 6814--6822. [Online]. Available: \url{https://doi.org/10.24963/ijcai.2023/763}

\bibitem{wen2015} M. Wen, R. Ehlers, and U. Topcu, ``Correct-by-synthesis reinforcement learning with temporal logic constraints,'' arXiv preprint arXiv:1503.01793, 2015. [Online]. Available: \url{https://arxiv.org/abs/1503.01793}

\bibitem{alshiekh2018} M. Alshiekh, R. Bloem, R. Ehlers, B. K\"{o}nighofer, S. Niekum, and U. Topcu, ``Safe reinforcement learning via shielding,'' in \emph{Proc. AAAI Conf. Artif. Intell.}, vol. 32, no. 1, 2018. [Online]. Available: \url{https://doi.org/10.1609/aaai.v32i1.11797}

\bibitem{jansen2020} N. Jansen, B. K\"{o}nighofer, S. Junges, A. C. Serban, and R. Bloem, ``Safe reinforcement learning via probabilistic shields,'' arXiv preprint arXiv:1807.06096, 2019. [Online]. Available: \url{https://arxiv.org/abs/1807.06096}

\bibitem{yang2023} W.-C. Yang, G. Marra, G. Rens, and L. De Raedt, ``Safe reinforcement learning via probabilistic logic shields,'' arXiv preprint arXiv:2303.03226, 2023. [Online]. Available: \url{https://arxiv.org/abs/2303.03226}

\bibitem{hamel2025} E. Hamel-De le Court, F. Belardinelli, and A. W. Goodall, ``Probabilistic shielding for safe reinforcement learning,'' arXiv preprint arXiv:2503.07671, 2025. [Online]. Available: \url{https://arxiv.org/abs/2503.07671}

\bibitem{hamel2026} E. Hamel-De le Court, T. Badings, A. Abate, F. Belardinelli, and F. Fabiano, ``Robust shielding for safe reinforcement learning,'' arXiv preprint arXiv:2606.00270, 2026. [Online]. Available: \url{https://arxiv.org/abs/2606.00270}

\bibitem{konighofer2025} B. K\"{o}nighofer, R. Bloem, N. Jansen, S. Junges, and S. Pranger, ``Shields for safe reinforcement learning,'' \emph{Communications of the ACM}, vol. 68, no. 11, pp. 80--90, 2025. [Online]. Available: \url{https://doi.org/10.1145/3715958}

\bibitem{pranger2026} S. Pranger and B. K\"{o}nighofer, ``Easy-to-use shielding for reinforcement learning,'' arXiv preprint arXiv:2606.03804, 2026. [Online]. Available: \url{https://arxiv.org/abs/2606.03804}

\bibitem{wang2025agentspec} H. Wang, C. M. Poskitt, and J. Sun, ``AgentSpec: Customizable runtime enforcement for safe and reliable LLM agents,'' in \emph{Proc. IEEE/ACM 48th Int. Conf. Softw. Eng. (ICSE)}, 2026, pp. 2938--2950. [Online]. Available: \url{https://doi.org/10.1145/3744916.3764546}

\bibitem{xiang2025} Z. Xiang \emph{et al.}, ``GuardAgent: Safeguard LLM agents via knowledge-enabled reasoning,'' arXiv preprint arXiv:2406.09187, 2025. [Online]. Available: \url{https://arxiv.org/abs/2406.09187}

\bibitem{chen2025shieldagent} Z. Chen, M. Kang, and B. Li, ``ShieldAgent: Shielding agents via verifiable safety policy reasoning,'' in \emph{Proc. 42nd Int. Conf. Mach. Learn. (ICML)}, PMLR vol. 267, 2025. [Online]. Available: \url{https://arxiv.org/abs/2503.22738}

\bibitem{hu2026} W. Hu, H. Jing, H. Shi, C. Fan, H. Li, and Y. Song, ``OmniCompliance-100K: A multi-domain, rule-grounded, real-world safety compliance dataset,'' arXiv preprint arXiv:2603.13933, 2026. [Online]. Available: \url{https://arxiv.org/abs/2603.13933}

\bibitem{levy2026} I. Levy \emph{et al.}, ``ST-WebAgentBench: A benchmark for evaluating safety and trustworthiness in web agents,'' in \emph{Proc. Int. Conf. Learn. Represent. (ICLR)}, 2026. [Online]. Available: \url{https://arxiv.org/abs/2410.06703}

\bibitem{mohammadi2025} M. Mohammadi, Y. Li, J. Lo, and W. Yip, ``Evaluation and benchmarking of LLM agents: A survey,'' arXiv preprint arXiv:2507.21504, 2025. [Online]. Available: \url{https://arxiv.org/abs/2507.21504}

\bibitem{turan2026} E. Turan, ``Oversight has a capacity: Calibrating agent guards to a subjective, fatiguing human,'' arXiv preprint arXiv:2606.08919, 2026. [Online]. Available: \url{https://arxiv.org/abs/2606.08919}

\bibitem{bondi2022} E. Bondi \emph{et al.}, ``Role of human-AI interaction in selective prediction,'' in \emph{Proc. AAAI Conf. Artif. Intell.}, vol. 36, no. 5, 2022, pp. 5286--5294. [Online]. Available: \url{https://doi.org/10.1609/aaai.v36i5.20465}

\end{thebibliography}

\end{document}
